% data-dict.tex
% Data Dictionary for MIMIC-III Bayesian Hierarchical Modeling Project
% PHP2530 Final Project - Brown University
% Author: Adam Kurth
% Last Updated: 2026-03-26

\documentclass[11pt]{article}
\usepackage[margin=1in]{geometry}
\usepackage{booktabs}
\usepackage{longtable}
\usepackage{array}
\usepackage{xcolor}
\usepackage{hyperref}
\usepackage{enumitem}
\usepackage{textcomp}
\usepackage{amsmath}
\usepackage{amssymb}

\definecolor{sectionblue}{RGB}{0,51,102}
\newcolumntype{L}[1]{>{\raggedright\arraybackslash}p{#1}}
\newcolumntype{C}[1]{>{\centering\arraybackslash}p{#1}}

\title{\textbf{Data Dictionary}\\[0.5em]
\large MIMIC-III Clinical Database (v1.4)\\
Bayesian Hierarchical Survival Modeling}
\author{PHP2530 Final Project}
\date{March 2026}

\begin{document}
\maketitle
\tableofcontents
\newpage

% ============================================================================
\section{Overview and Data Structure}
% ============================================================================

This data dictionary documents all variables extracted from the MIMIC-III Clinical Database for Bayesian hierarchical survival modeling. The preprocessing script (\texttt{data-prep.R}) creates intermediate datasets and \textbf{final analytical objects} for modeling.

\textbf{Primary Outputs} (loaded via \texttt{source("load.R")}):
\begin{itemize}[nosep]
    \item \texttt{X}: Static baseline predictor matrix ($N \times 66$)
    \item \texttt{y}: Binary ICU mortality outcome vector (length $N$)
    \item \texttt{S\_mat}: Time-varying predictor matrix ($\sum T_i \times Q$)
    \item \texttt{y\_S}: Discrete-time outcome vector (length $\sum T_i$)
\end{itemize}

See Section~\ref{sec:final-objects} for comprehensive documentation of these objects.

\subsection{Conceptual Framework}

For patient $i = 1, \ldots, N$ with $T_i$ discrete time periods, we define three \textbf{intermediate data frames} (for exploration) and two \textbf{final analytical objects} (for modeling):

\textbf{Intermediate Data Frames} (retain all variables, including IDs and outcomes):

\begin{enumerate}[leftmargin=*]
    \item \textbf{Static Baseline Covariate Matrix} $\mathbf{X}$ \texttt{(ready\_df)} $\longrightarrow$ cleaned to \texttt{X, y}
    \begin{itemize}[nosep]
        \item Dimension: $N \times P$ where $N$ = number of admissions, $P$ = number of covariates
        \item Each row represents ONE admission at ``Time Zero'' (ICU admission)
        \item Contains: demographics, aggregated vitals/labs, comorbidities, procedures, interventions
        \item These covariates are \textbf{constant within patient} --- they do not change over time
    \end{itemize}

    \item \textbf{Time-Varying Covariate Matrix} $\mathbf{Z}_{i,t}$ \texttt{(longitudinal\_df)}
    \begin{itemize}[nosep]
        \item Structure: One row per measurement occasion (patient $\times$ time point)
        \item $Z_{i,t}$ = covariate vector for patient $i$ at time $t$
        \item Contains: vital signs (HR, BP, SpO2, temp, RR) with precise timestamps
        \item These covariates \textbf{change over time} --- different values at each $t$
    \end{itemize}

    \item \textbf{Discrete-Time Survival Format} \texttt{(person\_period\_df)} $\longrightarrow$ cleaned to \texttt{S\_mat, y\_S}
    \begin{itemize}[nosep]
        \item Dimension: $\sum_{i=1}^{N} T_i$ rows (one row per patient per period at risk)
        \item Structure: $(i, t)$ indexed --- patient $i$, period $t$
        \item Contains: discrete outcome $y_{i,t}$, static covariates $X_i$, period-aggregated $Z_{i,t}$
        \item Designed for Polya-Gamma augmented Gibbs sampling
    \end{itemize}
\end{enumerate}

\subsection{Data Source Options}

The preprocessing script supports two data sources via \texttt{DATA\_SOURCE}:

\begin{itemize}[leftmargin=*, nosep]
    \item \texttt{DATA\_SOURCE <- "demo"}: MIMIC-III Demo ($\sim$100 patients) --- for development
    \item \texttt{DATA\_SOURCE <- "full"}: Complete MIMIC-III from PhysioNet ($\sim$46,000 patients)
\end{itemize}

\subsection{Output Files}

All files saved to \texttt{data/rds/} as R serialized objects:

\begin{longtable}{L{4.5cm}L{9.5cm}}
\toprule
\textbf{File} & \textbf{Description} \\
\midrule
\endhead
\multicolumn{2}{l}{\textit{Intermediate Data Frames (for exploration/debugging)}} \\
\texttt{ready\_df.rds} & Static baseline matrix before final exclusions \\
\texttt{longitudinal\_df.rds} & Time-varying measurements $\mathbf{Z}_{i,t}$ \\
\texttt{person\_period\_df.rds} & Discrete-time survival format (raw) \\
\midrule
\multicolumn{2}{l}{\textit{Final Analytical Objects (primary use --- loaded by \texttt{load.R})}} \\
\texttt{X\_y.rds} & List: \texttt{\{X, y\}} --- Static predictor matrix and outcome \\
\texttt{S\_y.rds} & List: \texttt{\{S, y\}} --- Time-varying matrix and discrete outcome \\
\bottomrule
\end{longtable}

\textbf{Important}: The final analytical objects (\texttt{X\_y.rds} and \texttt{S\_y.rds}) are the primary outputs for modeling. The intermediate data frames are retained for exploration but have NOT undergone final cleaning (exclusion of IDs, outcomes, multicollinear variables, etc.).

% ============================================================================
\section{Time Zero and Survival Framework}
% ============================================================================

\subsection{Time Zero Definition}

\textbf{Time Zero} is defined as the \textbf{first ICU admission} within each hospital stay (\texttt{icu\_intime}). This establishes a coherent survival framework:

\begin{center}
\begin{tabular}{ll}
\toprule
\textbf{Component} & \textbf{Definition} \\
\midrule
Time origin & First ICU admission timestamp (\texttt{icu\_intime}) \\
Event & Death during hospitalization (\texttt{event\_status = 1}) \\
Censoring & Discharge alive (\texttt{event\_status = 0}) \\
Time scale & Hours from ICU admission to death/discharge \\
\bottomrule
\end{tabular}
\end{center}

\textbf{Justification}: The ICU admission marks entry to highest-acuity care. Survival from this point is the clinically relevant outcome. Using the \textit{first} ICU stay avoids survivorship bias from conditioning on patients who survived long enough for subsequent transfers.

\subsection{Continuous-Time Outcome}

For the continuous-time framework (\texttt{ready\_df}):
\begin{itemize}[leftmargin=*, nosep]
    \item \texttt{time\_to\_event\_hours}: $T_i = $ hours from ICU admission to event/censoring
    \item \texttt{event\_status}: $\delta_i = 1$ if died, $\delta_i = 0$ if censored
\end{itemize}

\subsection{Discrete-Time Outcome}

For the discrete-time framework (\texttt{person\_period\_df}), continuous time is discretized into 24-hour periods:

\begin{itemize}[leftmargin=*, nosep]
    \item Period 1: hours 0--24, Period 2: hours 24--48, etc.
    \item Maximum 7 periods (capped for computational tractability)
    \item \texttt{y\_it}: binary outcome --- 1 if event in period $t$, 0 otherwise
\end{itemize}

\textbf{Example}:
\begin{center}
\begin{tabular}{lccccc}
\toprule
\textbf{Patient Type} & $y_{i,1}$ & $y_{i,2}$ & $y_{i,3}$ & $y_{i,4}$ & \textbf{Interpretation} \\
\midrule
Dies in period 3 & 0 & 0 & 1 & --- & Event in final period \\
Censored in period 4 & 0 & 0 & 0 & 0 & No event (discharged alive) \\
\bottomrule
\end{tabular}
\end{center}

% ============================================================================
\section{Final Analytical Objects (Primary Outputs)}
\label{sec:final-objects}
% ============================================================================

This section documents the \textbf{primary analytical objects} produced by \texttt{data-prep.R} and loaded by \texttt{load.R}. These are the clean, analysis-ready matrices and vectors for Bayesian modeling.

\subsection{Loading the Data}

Use \texttt{source("load.R")} to load all final objects into your R environment:

\begin{verbatim}
source("load.R")
# Objects now available:
#   X      - Static baseline predictor matrix (N x P)
#   y      - Binary outcome vector (length N)
#   S_mat  - Time-varying predictor matrix (sum of T_i x Q)
#   y_S    - Discrete-time outcome vector (length sum of T_i)
\end{verbatim}

\subsection{Object Relationship Diagram}

\begin{center}
\begin{tabular}{ccc}
\textbf{Static Analysis} & & \textbf{Discrete-Time Analysis} \\
\midrule
$\mathbf{X}$ (123 $\times$ 66) & $\longrightarrow$ & $\mathbf{S}$ ($\sum T_i$ $\times$ Q) \\
$\downarrow$ & & $\downarrow$ \\
$\mathbf{y}$ (123 elements) & & $\mathbf{y}_S$ ($\sum T_i$ elements) \\
\midrule
\textit{One row per admission} & & \textit{One row per person-period} \\
\textit{BART, logistic regression} & & \textit{Polya-Gamma Gibbs sampling} \\
\end{tabular}
\end{center}

\subsection{X: Static Baseline Predictor Matrix}

\textbf{Dimensions}: $N \times P$ where $N$ = number of admissions (after complete-case filtering), $P$ = number of predictors (66 columns in demo data).

\textbf{Source}: Derived from \texttt{ready\_df} after:
\begin{enumerate}[nosep]
    \item Exclusion of IDs, outcomes, timestamps, and redundant variables (Section \ref{sec:exclusions})
    \item Ethnicity recoding to numeric factors (Section \ref{sec:ethnicity})
    \item Complete-case analysis (\texttt{na.omit()})
    \item Removal of \texttt{hadm\_id} and \texttt{icu\_mortality} (extracted as \texttt{y})
\end{enumerate}

\textbf{Column Categories in X}:

\begin{longtable}{L{3.5cm}L{10.5cm}}
\toprule
\textbf{Category} & \textbf{Variables Included} \\
\midrule
\endhead
\multicolumn{2}{l}{\textit{Categorical Factors (Multi-level)}} \\
Demographics & admission\_type, insurance, ethnicity (0--7), gender, first\_careunit \\
Clinical Factors & age\_group (Adult/Elderly/Very Old), los\_group (Short/Prolonged/Extended), gcs\_status, complexity\_status, comorbidity\_level \\
Vital Sign Status & hr\_status, rr\_status, spo2\_status, temp\_status \\
Lab Status & bun\_status, creat\_status, glucose\_status, hgb\_status, plt\_status, k\_status, na\_status, wbc\_status \\
\midrule
\multicolumn{2}{l}{\textit{Binary Indicators (as Factors: 0/1)}} \\
Comorbidities & icd\_sepsis, icd\_heart\_failure, icd\_aki, icd\_ckd, icd\_resp\_failure, icd\_pneumonia, icd\_diabetes, icd\_liver\_disease, icd\_malignancy, icd\_stroke, icd\_mi, icd\_copd \\
Interventions & is\_ventilated, received\_dialysis, has\_central\_line, received\_transfusion, is\_surgical, on\_antibiotic, on\_anticoagulant, on\_insulin, on\_sedative \\
Microbiology & positive\_blood\_culture, any\_positive\_culture \\
Other & multiple\_icu\_stays, callout\_requested, hr\_unstable, temp\_unstable \\
\midrule
\multicolumn{2}{l}{\textit{Continuous (Retained)}} \\
Stay Characteristics & los\_icu, n\_icu\_stays \\
Counts & n\_diagnoses, n\_procedures, n\_prescriptions \\
Acuity & gcs\_mean\_val, vasopressor\_duration\_hrs, drg\_severity, drg\_mortality \\
\bottomrule
\end{longtable}

\subsection{y: Outcome Vector}

\textbf{Dimensions}: Length $N$ (same as number of rows in \texttt{X}).

\textbf{Definition}: \texttt{icu\_mortality} --- binary indicator of ICU mortality.
\begin{itemize}[nosep]
    \item \texttt{y = 1}: Patient died during ICU stay
    \item \texttt{y = 0}: Patient survived ICU stay (discharged alive or transferred)
\end{itemize}

\textbf{Relationship to other outcomes}:
\begin{itemize}[nosep]
    \item More conservative than \texttt{hospital\_expire\_flag} (hospital-wide mortality)
    \item Temporally bounded to the ICU admission window
    \item Used as the target for BART classification and static logistic regression
\end{itemize}

\subsection{S\_mat: Time-Varying Predictor Matrix}

\textbf{Dimensions}: $\left(\sum_{i=1}^{N} T_i\right) \times Q$ where $T_i$ = number of discrete periods patient $i$ is at risk, $Q$ = number of time-varying predictors.

\textbf{Source}: Derived from \texttt{person\_period\_df} after:
\begin{enumerate}[nosep]
    \item Filtering to matched cohort (\texttt{hadm\_id \%in\% valid\_hadm\_ids})
    \item Same variable exclusions as \texttt{X}
    \item Ethnicity recoding to match \texttt{X}
    \item Complete-case analysis (\texttt{na.omit()})
    \item Removal of \texttt{y\_it}, \texttt{icu\_mortality}, and \texttt{hadm\_id}
\end{enumerate}

\textbf{Key Variable}: \texttt{period} (integer 1--7) is \textbf{retained} in \texttt{S\_mat} to support baseline hazard approximations in discrete-time models.

\textbf{Structure}: Each patient contributes $T_i$ rows (one per period at risk). Static covariates are repeated across periods; time-varying vitals are period-specific aggregates (mean, min, max per 24-hour window).

\textbf{Example rows for one patient}:
\begin{center}
\begin{tabular}{cccccc}
\toprule
\textbf{row} & \textbf{period} & \textbf{age\_group} & \textbf{hr\_mean} & \textbf{temp\_mean} & \textbf{...} \\
\midrule
1 & 1 & Elderly & 85.2 & 37.1 & ... \\
2 & 2 & Elderly & 88.7 & 37.3 & ... \\
3 & 3 & Elderly & 92.1 & 37.8 & ... \\
\bottomrule
\end{tabular}
\end{center}

\subsection{y\_S: Discrete-Time Outcome Vector}

\textbf{Dimensions}: Length $\sum_{i=1}^{N} T_i$ (same as number of rows in \texttt{S\_mat}).

\textbf{Definition}: \texttt{y\_it} --- binary indicator of event in period $t$ for patient $i$.
\begin{itemize}[nosep]
    \item \texttt{y\_S[row] = 1}: Event (death) occurred in this period
    \item \texttt{y\_S[row] = 0}: No event in this period (survived to next period or censored)
\end{itemize}

\textbf{Key property}: For each patient, \texttt{y\_S = 1} in \textbf{at most one row} (the final period if they died). All prior periods have \texttt{y\_S = 0}.

\textbf{Use case}: Polya-Gamma augmented Gibbs sampling for discrete-time survival models:
\[
\text{logit}\, P(y_{it} = 1 \mid y_{i,t-1} = 0) = \alpha_t + S_{it}' \beta
\]

\subsection{Correspondence Between Objects}

\begin{center}
\begin{tabular}{lll}
\toprule
\textbf{Static Object} & \textbf{Time-Varying Object} & \textbf{Relationship} \\
\midrule
\texttt{X} & \texttt{S\_mat} & Static vars repeated; S adds \texttt{period} + time-varying vitals \\
\texttt{y} & \texttt{y\_S} & \texttt{y} = patient-level; \texttt{y\_S} = period-level (sums to same \# events) \\
\texttt{hadm\_id} (row names) & \texttt{hadm\_id} (in S before drop) & Links patients across matrices \\
\bottomrule
\end{tabular}
\end{center}

% ============================================================================
\section{Static Baseline Matrix X (ready\_df)}
% ============================================================================

This dataset contains \textbf{one row per hospital admission}. All variables represent information available at or aggregated across the ICU stay, treated as ``baseline'' for modeling purposes.

\subsection{Identification Variables}

\begin{longtable}{L{3cm}L{2cm}L{9cm}}
\toprule
\textbf{Variable} & \textbf{Type} & \textbf{Description} \\
\midrule
\endhead
subject\_id & Factor & Unique patient identifier (for patient-level random effects) \\
hadm\_id & Factor & Unique hospital admission identifier (unit of analysis) \\
icustay\_id & Integer & First ICU stay identifier within admission \\
\bottomrule
\end{longtable}

\textbf{Note}: All identification variables are \textbf{excluded} from the final predictor matrix \texttt{X} in Section 17. They are retained in \texttt{ready\_df} for reference and hierarchical modeling.

\subsection{Demographics}

\begin{longtable}{L{3cm}L{2cm}L{9cm}}
\toprule
\textbf{Variable} & \textbf{Type} & \textbf{Description} \\
\midrule
\endhead
age & Numeric & Age at admission (years), capped at 90 per MIMIC convention \\
age\_group & Factor & Discretized: Adult ($<$65), Elderly ($\geq$65), Very Old ($\geq$80) \\
gender & Factor & Patient sex (M/F) \\
ethnicity & Factor & Recoded to numeric levels 0--7 (see Section 5.8) \\
insurance & Factor & Primary insurance type \\
\bottomrule
\end{longtable}

\textbf{Variables removed}:
\begin{itemize}[leftmargin=*, nosep]
    \item marital\_status --- dropped in final cleaning
    \item scaled\_age --- dropped (multicollinear with age)
\end{itemize}

\subsection{Admission \& ICU Variables}

\begin{longtable}{L{3.5cm}L{2cm}L{8.5cm}}
\toprule
\textbf{Variable} & \textbf{Type} & \textbf{Description} \\
\midrule
\endhead
admission\_type & Factor & ELECTIVE, EMERGENCY, URGENT \\
first\_careunit & Factor & ICU type: MICU, SICU, CCU, CSRU, TSICU \\
first\_service & Factor & First hospital service \\
los\_icu & Numeric & First ICU length of stay (days) \\
los\_group & Factor & Discretized: Short ($<$3d), Prolonged ($\geq$3d), Extended ($\geq$7d) \\
n\_icu\_stays & Integer & Total ICU stays during admission \\
multiple\_icu\_stays & Binary Factor & 1 if n\_icu\_stays $>1$ \\
is\_surgical & Binary & 1 if any surgical service involvement \\
n\_service\_changes & Integer & Number of service transfers \\
n\_transfers & Integer & Number of unit transfers \\
\bottomrule
\end{longtable}

\textbf{Variables removed}:
\begin{itemize}[leftmargin=*, nosep]
    \item admittime, dischtime, deathtime, icu\_intime, icu\_outtime --- timestamps excluded
    \item scaled\_los\_icu --- multicollinear with los\_icu
    \item total\_icu\_los --- redundant with los\_icu
    \item n\_unique\_units --- dropped in final cleaning
\end{itemize}

\subsection{Survival Outcomes}

\begin{longtable}{L{4cm}L{2cm}L{8cm}}
\toprule
\textbf{Variable} & \textbf{Type} & \textbf{Description} \\
\midrule
\endhead
icu\_mortality & Binary & \textbf{TARGET}: 1 = died during ICU stay \\
\bottomrule
\end{longtable}

\textbf{Variables excluded from X (leakage/outcomes)}:
\begin{itemize}[leftmargin=*, nosep]
    \item event\_status --- hospital mortality indicator (1 = died during hospitalization)
    \item hospital\_expire\_flag --- redundant with event\_status
    \item time\_to\_event\_hours, time\_to\_event\_days --- time-to-event (use for survival modeling only)
    \item end\_time --- death/discharge timestamp
\end{itemize}

\textbf{Note}: The target variable \texttt{icu\_mortality} is a stricter definition than \texttt{event\_status}:
\begin{itemize}[leftmargin=*, nosep]
    \item \texttt{icu\_mortality = 1}: Patient died \textbf{during ICU stay} (deathtime $\leq$ icu\_outtime)
    \item \texttt{event\_status = 1}: Patient died during hospitalization (may have been after ICU discharge)
\end{itemize}

\subsection{Vital Sign Summaries}

Per-admission aggregates from CHARTEVENTS. \textbf{Note}: Only mean and SD statistics are retained; min/max values are dropped due to multicollinearity with mean values.

\begin{center}
\texttt{hr\_mean\_val, hr\_sd\_val} (then removed after factorization)
\end{center}

\begin{longtable}{L{3cm}L{2.5cm}L{2.5cm}L{5.5cm}}
\toprule
\textbf{Vital Sign} & \textbf{CareVue ID} & \textbf{MetaVision ID} & \textbf{Units} \\
\midrule
\endhead
hr (heart rate) & 211 & 220045 & bpm \\
resp\_rate & 618 & 220210 & breaths/min \\
spo2 & 646 & 220277 & \% \\
temp\_c & 676, 678$\to$C & 223762 & \textdegree C \\
gcs & 198 & 220739 & 3--15 \\
\bottomrule
\end{longtable}

\textbf{Variables dropped due to high missingness} ($>25\%$ in demo dataset):
\begin{itemize}[leftmargin=*, nosep]
    \item sbp (systolic BP) --- $\sim$57\% missing
    \item dbp (diastolic BP) --- $\sim$81\% missing
    \item map (mean arterial pressure) --- $\sim$58\% missing
\end{itemize}

\textbf{Statistics computed} (intermediate, later dropped after factorization):
\begin{itemize}[leftmargin=*, nosep]
    \item \texttt{mean\_val}: Average across all measurements $\to$ used to create categorical factors
    \item \texttt{sd\_val}: Variability (instability indicator) $\to$ used for \_unstable flags
    \item \texttt{min\_val, max\_val, n\_obs}: \textbf{Dropped} due to multicollinearity
\end{itemize}

\subsection{Laboratory Summaries}

Per-admission aggregates from LABEVENTS. \textbf{Note}: Only mean values are retained and converted to categorical factors; min/max/first/last values are dropped.

\begin{longtable}{L{3cm}L{1.5cm}L{2cm}L{6.5cm}}
\toprule
\textbf{Lab Test} & \textbf{ID} & \textbf{Units} & \textbf{Clinical Relevance} \\
\midrule
\endhead
creatinine & 50912 & mg/dL & Renal function (SOFA) \\
platelets & 51265 & K/uL & Coagulation (SOFA) \\
wbc & 51301 & K/uL & Infection/immune \\
hemoglobin & 51222 & g/dL & Anemia, bleeding \\
glucose & 50931 & mg/dL & Metabolic status \\
bun & 51006 & mg/dL & Renal function \\
sodium & 50983 & mEq/L & Electrolyte balance \\
potassium & 50971 & mEq/L & Cardiac/renal \\
\bottomrule
\end{longtable}

\textbf{Labs dropped due to high missingness} ($>25\%$ in demo dataset):
\begin{itemize}[leftmargin=*, nosep]
    \item lactate ($\sim$21\% missing)
    \item bilirubin ($\sim$32\% missing)
    \item pao2, pco2 ($\sim$29\% missing)
    \item ph ($\sim$26\% missing)
\end{itemize}

\textbf{Statistics computed} (intermediate, later dropped after factorization):
\begin{itemize}[leftmargin=*, nosep]
    \item \texttt{mean\_val}: Average across measurements $\to$ used to create categorical factors
    \item \texttt{first\_val, last\_val, min\_val, max\_val, n\_labs}: \textbf{Dropped} due to multicollinearity
\end{itemize}

\subsection{Comorbidity Flags (ICD-9)}

\begin{longtable}{L{3.5cm}L{1.5cm}L{3cm}L{5cm}}
\toprule
\textbf{Variable} & \textbf{Type} & \textbf{ICD-9 Pattern} & \textbf{Description} \\
\midrule
\endhead
n\_diagnoses & Integer & --- & Total diagnosis codes \\
icd\_sepsis & Binary Factor & 995.91, 995.92, 038.x & Sepsis/septicemia \\
icd\_heart\_failure & Binary Factor & 428.x & CHF \\
icd\_aki & Binary Factor & 584.x & Acute kidney injury \\
icd\_ckd & Binary Factor & 585.x & Chronic kidney disease \\
icd\_resp\_failure & Binary Factor & 518.8x & Respiratory failure \\
icd\_pneumonia & Binary Factor & 480--486 & Pneumonia \\
icd\_diabetes & Binary Factor & 250.x & Diabetes mellitus \\
icd\_liver\_disease & Binary Factor & 571.x & Chronic liver disease \\
icd\_malignancy & Binary Factor & 140--208 & Cancer \\
icd\_stroke & Binary Factor & 430--438 & Cerebrovascular disease \\
icd\_mi & Binary Factor & 410.x & Myocardial infarction \\
icd\_copd & Binary Factor & 490--496 & COPD \\
comorbidity\_level & Factor & --- & Low/Moderate ($<$3), High ($\geq$3), Very High ($\geq$5) \\
\bottomrule
\end{longtable}

\textbf{Variables removed}:
\begin{itemize}[leftmargin=*, nosep]
    \item comorbidity\_count --- exact sum of icd\_* flags (redundant)
    \item scaled\_comorbidity --- multicollinear
    \item scaled\_n\_diagnoses --- multicollinear with n\_diagnoses
\end{itemize}

\subsection{Procedures \& Interventions}

\begin{longtable}{L{4.5cm}L{2cm}L{7.5cm}}
\toprule
\textbf{Variable} & \textbf{Type} & \textbf{Description} \\
\midrule
\endhead
n\_procedures & Integer & Total ICD-9 procedure codes \\
is\_ventilated & Binary Factor & Mechanical ventilation (96.7x) \\
received\_dialysis & Binary Factor & Hemodialysis (39.95) \\
has\_central\_line & Binary Factor & Central venous catheter \\
received\_transfusion & Binary Factor & Blood transfusion (99.0x) \\
n\_vasopressor\_doses & Integer & Total vasopressor events \\
vasopressor\_duration\_hrs & Numeric & Total hours on vasopressors \\
complexity\_status & Factor & Standard vs. High Complexity \\
\bottomrule
\end{longtable}

\textbf{Variables removed}:
\begin{itemize}[leftmargin=*, nosep]
    \item received\_vasopressors --- redundant with vasopressor\_duration\_hrs ($>0$)
    \item scaled\_n\_procedures --- multicollinear with n\_procedures
\end{itemize}

\subsection{DRG Severity Scores}

\begin{longtable}{L{4cm}L{2cm}L{8cm}}
\toprule
\textbf{Variable} & \textbf{Type} & \textbf{Description} \\
\midrule
\endhead
drg\_mortality\_score & Integer & DRG mortality risk (1--4, higher = worse) \\
drg\_severity\_score & Integer & DRG severity (1--4, higher = more severe) \\
\bottomrule
\end{longtable}

% ============================================================================
\section{Time-Varying Matrix Z (longitudinal\_df)}
% ============================================================================

This dataset contains \textbf{one row per vital sign measurement occasion}. Structure: $Z_{i,t}$ = covariate vector for patient $i$ at time $t$.

\subsection{Key Variables}

\begin{longtable}{L{4cm}L{2cm}L{8cm}}
\toprule
\textbf{Variable} & \textbf{Type} & \textbf{Description} \\
\midrule
\endhead
subject\_id & Factor & Patient identifier \\
hadm\_id & Factor & Admission identifier \\
icustay\_id & Integer & ICU stay identifier \\
charttime & POSIXct & Measurement timestamp \\
hours\_since\_icu\_admit & Numeric & $t$ = hours since ICU admission (time index) \\
scaled\_time & Numeric & Standardized time \\
\bottomrule
\end{longtable}

\subsection{Vital Signs (Raw Values)}

\begin{longtable}{L{3cm}L{2cm}L{9cm}}
\toprule
\textbf{Variable} & \textbf{Type} & \textbf{Description} \\
\midrule
\endhead
hr & Numeric & Heart rate (bpm) \\
sbp & Numeric & Systolic blood pressure (mmHg) \\
dbp & Numeric & Diastolic blood pressure (mmHg) \\
map & Numeric & Mean arterial pressure (mmHg) \\
resp\_rate & Numeric & Respiratory rate (breaths/min) \\
spo2 & Numeric & Oxygen saturation (\%) \\
temp\_c & Numeric & Temperature (\textdegree C) \\
gcs & Numeric & Glasgow Coma Scale (3--15) \\
\bottomrule
\end{longtable}

\subsection{Scaled Vitals}

\begin{center}
\texttt{scaled\_hr, scaled\_sbp, scaled\_dbp, scaled\_map, scaled\_resp\_rate, scaled\_spo2, scaled\_temp\_c}
\end{center}

\subsection{Outcome Information (Merged from Cohort)}

\begin{longtable}{L{4cm}L{2cm}L{8cm}}
\toprule
\textbf{Variable} & \textbf{Type} & \textbf{Description} \\
\midrule
\endhead
event\_status & Binary & 1 = died, 0 = censored \\
time\_to\_event\_hours & Numeric & Patient's total survival time \\
end\_time & POSIXct & Death or discharge timestamp \\
first\_careunit & Factor & ICU type \\
age & Numeric & Age at admission \\
gender & Factor & Patient sex \\
\bottomrule
\end{longtable}

\subsection{Filtering Applied}

The dataset is restricted to:
\begin{itemize}[leftmargin=*, nosep]
    \item Measurements after ICU admission (\texttt{hours\_since\_icu\_admit} $\geq$ 0)
    \item First 7 days post-ICU admission (\texttt{hours\_since\_icu\_admit} $\leq$ 168)
\end{itemize}

% ============================================================================
\section{Discrete-Time Format (person\_period\_df)}
% ============================================================================

This dataset contains \textbf{one row per patient per discrete time period}. Designed for discrete-time survival modeling with Polya-Gamma augmentation.

\subsection{Why Discrete-Time Survival?}

Continuous-time survival models (Cox, Weibull) can be computationally intensive in Bayesian settings. Discrete-time models offer advantages:

\begin{enumerate}[leftmargin=*, nosep]
    \item \textbf{Logistic regression formulation}: Transforms survival into sequence of binary outcomes
    \item \textbf{Polya-Gamma augmentation}: Enables conjugate Gibbs sampling (no Metropolis-Hastings)
    \item \textbf{Natural for time-varying covariates}: One covariate value per discrete period
    \item \textbf{Computational efficiency}: Essential for large datasets like MIMIC-III
\end{enumerate}

\subsection{Mathematical Formulation}

The discrete-time hazard:
\[
h_{it} = \Pr(T_i \in (a_{t-1}, a_t] \mid T_i > a_{t-1}, X_i, Z_{i,t})
\]

Logistic regression model:
\[
\text{logit}(h_{it}) = \alpha_t + X_i'\beta + Z_{i,t}'\gamma + u_i
\]

where:
\begin{itemize}[leftmargin=*, nosep]
    \item $\alpha_t$: period-specific baseline hazard (intercept)
    \item $X_i$: static covariates (repeated across periods)
    \item $Z_{i,t}$: time-varying covariates (different each period)
    \item $u_i \sim N(0, \sigma^2_u)$: patient-level random effect (frailty)
\end{itemize}

\subsection{Polya-Gamma Augmentation}

The logistic likelihood is not conjugate to Gaussian priors. Polya-Gamma augmentation (Polson et al., 2013) introduces auxiliary variables $\omega_{it} \sim \text{PG}(1, \eta_{it})$ that make the conditional posterior of $\beta$ Gaussian:

\[
\beta \mid \omega, y \sim \text{Normal}(\cdot, \cdot)
\]

This enables \textbf{fully conjugate Gibbs sampling} --- dramatically improving computational efficiency and mixing for large datasets.

\subsection{Key Variables}

\begin{longtable}{L{4cm}L{2cm}L{8cm}}
\toprule
\textbf{Variable} & \textbf{Type} & \textbf{Description} \\
\midrule
\endhead
hadm\_id & Factor & Admission identifier \\
subject\_id & Factor & Patient identifier \\
period & Integer & Time period: 1, 2, \ldots, 7 (days) \\
max\_period & Integer & Total periods at risk for this patient \\
y\_it & Binary & \textbf{Discrete outcome}: 1 if event in period $t$ \\
event\_status & Binary & Original continuous-time indicator \\
\bottomrule
\end{longtable}

\subsection{Static Covariates (from X)}

These are the $X_i$ variables \textbf{repeated across all periods} for the same patient:

\begin{longtable}{L{4cm}L{2cm}L{8cm}}
\toprule
\textbf{Variable} & \textbf{Type} & \textbf{Description} \\
\midrule
\endhead
age & Numeric & Age at admission \\
gender & Factor & Patient sex \\
first\_careunit & Factor & ICU type \\
comorbidity\_count & Integer & Number of comorbidities \\
received\_vasopressors & Binary & Any vasopressor use \\
is\_ventilated & Binary & Mechanical ventilation \\
icd\_sepsis & Binary & Sepsis diagnosis \\
drg\_mortality\_score & Integer & DRG mortality risk \\
los\_icu & Numeric & ICU length of stay \\
\bottomrule
\end{longtable}

\subsection{Time-Varying Covariates (from Z)}

Period-level aggregations of vital signs --- these are the $Z_{i,t}$ that \textbf{change across periods}:

\begin{longtable}{L{4cm}L{2cm}L{8cm}}
\toprule
\textbf{Variable} & \textbf{Type} & \textbf{Description} \\
\midrule
\endhead
hr\_mean & Numeric & Mean heart rate in period $t$ \\
hr\_max & Numeric & Maximum heart rate in period $t$ \\
hr\_n & Integer & Number of HR measurements in period $t$ \\
map\_mean & Numeric & Mean MAP in period $t$ \\
resp\_rate\_mean & Numeric & Mean respiratory rate in period $t$ \\
spo2\_mean & Numeric & Mean SpO2 in period $t$ \\
temp\_c\_mean & Numeric & Mean temperature in period $t$ \\
\bottomrule
\end{longtable}

\subsection{Missingness Indicators}

Missing vitals in ICU settings may be informative (e.g., stable patients monitored less frequently). We create indicators:

\begin{longtable}{L{4cm}L{2cm}L{8cm}}
\toprule
\textbf{Variable} & \textbf{Type} & \textbf{Description} \\
\midrule
\endhead
hr\_missing & Binary & 1 if no HR measurements in period (before imputation) \\
temp\_c\_missing & Binary & 1 if no temperature measurements in period \\
map\_missing & Binary & 1 if no MAP measurements in period \\
\bottomrule
\end{longtable}

\subsection{Missing Data Imputation}

For time-varying covariates in person-period format:

\begin{enumerate}[leftmargin=*, nosep]
    \item \textbf{Create missingness indicators} (before imputation) --- allows model to learn if missingness is prognostic
    \item \textbf{LOCF (Last Observation Carried Forward)}: Within each patient, fill missing periods with most recent available value
    \item \textbf{Population median}: For patients with no prior measurements, use population median
\end{enumerate}

\textbf{Note}: LOCF artificially shrinks variance. Consider fully Bayesian imputation (treating missing values as latent parameters) for sensitivity analysis.

% ============================================================================
\section{Discretized Variables for BART Modeling}
% ============================================================================

This section documents the discretization of continuous variables into \textbf{multi-level categorical factors} for BART modeling. The implementation uses \texttt{factor(case\_when(...), levels = ...)} to create ordered factor variables with clinically meaningful categories.

\subsection{Rationale for Multi-Level Factorization}

The approach creates ordered categorical factors with 2--3 levels based on clinical thresholds:

\begin{enumerate}[leftmargin=*, nosep]
    \item \textbf{Clinical interpretability}: Categories match clinical severity grades (Normal $\to$ Mild $\to$ Severe)
    \item \textbf{SOFA/SIRS alignment}: Thresholds correspond to established scoring systems
    \item \textbf{Ordered structure}: Factor levels are ordered from normal to most severe
    \item \textbf{BART efficiency}: Multi-level factors enable efficient tree splits at multiple thresholds
    \item \textbf{Information preservation}: Distinguishes mild from severe abnormalities (vs.\ binary flags)
\end{enumerate}

\textbf{Pattern used}: \texttt{factor(case\_when(...), levels = c("Normal", "Mild", "Severe"))}

\subsection{Variables Removed in Section 14b}

The following high-cardinality and problematic variables are removed from \texttt{ready\_df}:

\begin{longtable}{L{3.5cm}L{10.5cm}}
\toprule
\textbf{Variable} & \textbf{Reason for Removal} \\
\midrule
\endhead
admission\_location & High-cardinality text (7+ categories with long strings) \\
discharge\_location & High-cardinality text (15+ categories); \textbf{LEAKAGE} (known only at discharge) \\
language & High-cardinality character (many languages + empty strings) \\
religion & High-cardinality character (many religions + empty strings) \\
diagnosis & Free-text primary diagnosis --- extremely high cardinality \\
dob & POSIXct datetime --- use \texttt{age} instead \\
dod & POSIXct datetime --- use \texttt{event\_status} instead \\
expire\_flag & Outcome indicator --- \textbf{LEAKAGE} (redundant with event\_status) \\
has\_chartevents\_data & Binary indicator --- all ICU patients have chart data (no variance) \\
\bottomrule
\end{longtable}

\subsection{Variables Removed Due to High Missingness (Section 14)}

These variables are dropped because they exceed 25\% missing in the demo dataset:

\begin{longtable}{L{4cm}L{3cm}L{7cm}}
\toprule
\textbf{Variable Pattern} & \textbf{Missingness} & \textbf{Notes} \\
\midrule
\endhead
dbp\_* & $\sim$81\% & Diastolic BP aggregates \\
map\_* & $\sim$58\% & Mean arterial pressure aggregates \\
sbp\_* & $\sim$57\% & Systolic BP aggregates \\
total\_input\_*, n\_input\_* & $\sim$45\% & Input volume data \\
n\_positive\_spec*, n\_unique\_org* & $\sim$44\% & Microbiology details \\
dod\_hosp & $\sim$37\% & Date of death (hospital) \\
bilirubin\_* & $\sim$32\% & Bilirubin labs \\
pao2\_*, pco2\_* & $\sim$29\% & Arterial blood gases \\
edregtime, edouttime & $\sim$28\% & ED timestamps \\
n\_blood\_cultures, n\_positive\_cult*, n\_organisms & $\sim$27\% & Culture counts \\
ph\_* & $\sim$26\% & pH labs \\
lactate\_* & $\sim$21\% & Lactate labs \\
dod\_ssn & $\sim$19\% & SSN death date \\
\bottomrule
\end{longtable}

\textbf{Note}: These thresholds are calibrated for the demo dataset. For the full MIMIC-III dataset, consider retaining variables with lower missingness.

\subsection{Multi-Level Demographic Factors}

\begin{longtable}{L{2.5cm}L{2.5cm}L{4cm}L{5cm}}
\toprule
\textbf{Variable} & \textbf{Type} & \textbf{Levels (Ordered)} & \textbf{Clinical Thresholds} \\
\midrule
\endhead
age\_group & Factor & Adult, Elderly, Very Old & $<65$, $\geq65$, $\geq80$ years \\
los\_group & Factor & Short, Prolonged, Extended & $<3$, $\geq3$, $\geq7$ days \\
multiple\_icu\_stays & Binary Factor & 0, 1 & n\_icu\_stays $>1$ \\
\bottomrule
\end{longtable}

\subsection{Multi-Level Vital Sign Factors}

All thresholds based on SIRS criteria and standard clinical cutoffs. Each variable has 2--3 ordered levels.

\begin{longtable}{L{2.5cm}L{5cm}L{6.5cm}}
\toprule
\textbf{Variable} & \textbf{Levels (Ordered)} & \textbf{Clinical Thresholds} \\
\midrule
\endhead
\multicolumn{3}{l}{\textbf{Glasgow Coma Scale}} \\
gcs\_status & Normal, Impaired, Severe Impairment & GCS $=15$, $<15$, $\leq8$ \\
\midrule
\multicolumn{3}{l}{\textbf{Heart Rate}} \\
hr\_status & Normal, Bradycardia, Tachycardia & 60--100, $<60$, $>100$ bpm \\
hr\_unstable & 0, 1 (binary factor) & HR SD $>15$ \\
\midrule
\multicolumn{3}{l}{\textbf{Respiratory Rate}} \\
rr\_status & Normal, Tachypnea, Severe Tachypnea & $\leq20$, $>20$, $>30$ /min \\
\midrule
\multicolumn{3}{l}{\textbf{Oxygen Saturation}} \\
spo2\_status & Normal, Hypoxia, Severe Hypoxia & $\geq95\%$, $<95\%$, $<90\%$ \\
\midrule
\multicolumn{3}{l}{\textbf{Temperature}} \\
temp\_status & Normal, Hypothermia, Fever & 36--38\textdegree C, $<36$\textdegree C, $>38$\textdegree C \\
temp\_unstable & 0, 1 (binary factor) & Temp SD $>0.5$ \\
\bottomrule
\end{longtable}

\subsection{Multi-Level Laboratory Factors}

Thresholds based on SOFA criteria and standard reference ranges.

\begin{longtable}{L{2.5cm}L{5cm}L{6.5cm}}
\toprule
\textbf{Variable} & \textbf{Levels (Ordered)} & \textbf{Clinical Thresholds} \\
\midrule
\endhead
\multicolumn{3}{l}{\textbf{Renal Function}} \\
bun\_status & Normal, High, Very High & $\leq20$, $>20$, $>40$ mg/dL \\
creat\_status & Normal, Elevated, Dysfunction & $<1.2$, $\geq1.2$, $\geq2.0$ mg/dL \\
\midrule
\multicolumn{3}{l}{\textbf{Glucose}} \\
glucose\_status & Normal, Hypoglycemia, Hyperglycemia & 70--180, $<70$, $>180$ mg/dL \\
\midrule
\multicolumn{3}{l}{\textbf{Hemoglobin}} \\
hgb\_status & Normal, Anemia, Severe Anemia & $\geq10$, $<10$, $<7$ g/dL \\
\midrule
\multicolumn{3}{l}{\textbf{Platelets}} \\
plt\_status & Normal, Thrombocytopenia, Severe Thrombocytopenia & $\geq150$, $<150$, $<50$ K/uL \\
\midrule
\multicolumn{3}{l}{\textbf{Electrolytes}} \\
k\_status & Normal, Hypokalemia, Hyperkalemia & 3.5--5.0, $<3.5$, $>5.0$ mEq/L \\
na\_status & Normal, Hyponatremia, Hypernatremia & 135--145, $<135$, $>145$ mEq/L \\
\midrule
\multicolumn{3}{l}{\textbf{White Blood Cells}} \\
wbc\_status & Normal, Leukopenia, Leukocytosis & 4--12, $<4$, $>12$ K/uL \\
\bottomrule
\end{longtable}

\subsection{Multi-Level Complexity Factors}

\begin{longtable}{L{3.5cm}L{4.5cm}L{6cm}}
\toprule
\textbf{Variable} & \textbf{Levels (Ordered)} & \textbf{Clinical Thresholds} \\
\midrule
\endhead
comorbidity\_level & Low/Moderate, High, Very High & $<3$, $\geq3$, $\geq5$ comorbidities \\
complexity\_status & Standard, High Complexity & n\_rx $>20$ OR n\_dx $>10$ OR n\_proc $>5$ \\
\bottomrule
\end{longtable}

\subsection{Binary Indicators (Converted to Factors)}

The following 0/1 integer variables are converted to factors using \texttt{as.factor()}:

\begin{itemize}[leftmargin=*, nosep]
    \item \textbf{ICD Comorbidities}: icd\_sepsis, icd\_heart\_failure, icd\_aki, icd\_ckd, icd\_resp\_failure, icd\_pneumonia, icd\_diabetes, icd\_liver\_disease, icd\_malignancy, icd\_stroke, icd\_mi, icd\_copd
    \item \textbf{Interventions}: is\_ventilated, received\_dialysis, has\_central\_line, received\_transfusion, is\_surgical
    \item \textbf{Medications}: on\_antibiotic, on\_anticoagulant, on\_insulin, on\_sedative
    \item \textbf{Cultures}: positive\_blood\_culture, any\_positive\_culture
    \item \textbf{Other}: callout\_requested
\end{itemize}

\subsection{Categorical Variables (Multi-level Factors)}

These remain as multi-level factors with their original levels:

\begin{longtable}{L{3cm}L{11cm}}
\toprule
\textbf{Variable} & \textbf{Levels} \\
\midrule
\endhead
admission\_type & ELECTIVE, EMERGENCY, URGENT \\
insurance & Medicare, Medicaid, Private, Government, Self Pay \\
first\_careunit & MICU, SICU, CCU, CSRU, TSICU \\
first\_service & MED, SURG, CMED, CSURG, NSURG, etc. \\
gender & M, F \\
\bottomrule
\end{longtable}

\subsection{Ethnicity Recoding}
\label{sec:ethnicity}

Ethnicity is recoded from text strings to numeric factor codes for consistency:

\begin{longtable}{L{6cm}L{2cm}L{6cm}}
\toprule
\textbf{Original Value} & \textbf{Code} & \textbf{Category} \\
\midrule
\endhead
WHITE & 1 & White \\
BLACK/AFRICAN AMERICAN & 2 & Black \\
ASIAN & 3 & Asian \\
AMERICAN INDIAN/ALASKA NATIVE... & 4 & American Indian/Alaska Native \\
HISPANIC OR LATINO, HISPANIC/LATINO - PUERTO RICAN & 5 & Hispanic/Latino \\
OTHER & 6 & Other \\
NATIVE HAWAIIAN OR OTHER PACIFIC ISLANDER & 7 & Pacific Islander \\
UNKNOWN/NOT SPECIFIED, OTHER/UNKNOWN, UNABLE TO OBTAIN, NA & 0 & Unknown \\
\bottomrule
\end{longtable}

Final factor levels: \texttt{c("0", "1", "2", "3", "4", "5", "6", "7")}

\subsection{Variables Removed After Factorization}

After creating the categorical factors, the \textbf{original continuous values} are removed to avoid multicollinearity:

\begin{itemize}[leftmargin=*, nosep]
    \item \textbf{Original vitals}: gcs\_sd\_val, gcs\_n\_obs, hr\_mean\_val, hr\_sd\_val, hr\_n\_obs, resp\_rate\_mean\_val, resp\_rate\_sd\_val, resp\_rate\_n\_obs, spo2\_mean\_val, spo2\_sd\_val, spo2\_n\_obs, temp\_c\_mean\_val, temp\_c\_sd\_val, temp\_c\_n\_obs
    \item \textbf{Original labs}: bun\_mean\_val, bun\_n\_labs, creatinine\_mean\_val, creatinine\_n\_labs, glucose\_mean\_val, glucose\_n\_labs, hemoglobin\_mean\_val, hemoglobin\_n\_labs, platelets\_mean\_val, platelets\_n\_labs, potassium\_mean\_val, potassium\_n\_labs, sodium\_mean\_val, sodium\_n\_labs, wbc\_mean\_val, wbc\_n\_labs
    \item \textbf{Scaled variables}: All variables starting with \texttt{scaled\_}
\end{itemize}

% ============================================================================
\section{Final Cleaning and X/y Extraction (Section 17)}
% ============================================================================

Section 17 of \texttt{data-prep.R} performs final cleaning to produce analysis-ready matrices for BART modeling.

\subsection{Excluded Variables}

The following variables are explicitly excluded from the final predictor matrix \texttt{X}:

\begin{longtable}{L{4cm}L{10cm}}
\toprule
\textbf{Variable Category} & \textbf{Variables Excluded} \\
\midrule
\endhead
\multicolumn{2}{l}{\textbf{Identification Variables}} \\
& subject\_id, hadm\_id, icustay\_id \\
\midrule
\multicolumn{2}{l}{\textbf{Outcome Variables (Leakage Prevention)}} \\
& event\_status, hospital\_expire\_flag, time\_to\_event\_hours, time\_to\_event\_days \\
\midrule
\multicolumn{2}{l}{\textbf{Timestamp Variables}} \\
& admittime, dischtime, deathtime, icu\_intime, icu\_outtime, end\_time \\
\midrule
\multicolumn{2}{l}{\textbf{High-Cardinality/Missing Variables}} \\
& discharge\_location, expire\_flag, dob, dod, dod\_hosp, dod\_ssn, diagnosis, marital\_status, religion, language \\
\midrule
\multicolumn{2}{l}{\textbf{Exact Multicollinearity (Scaled Duplicates)}} \\
& scaled\_age, scaled\_los\_icu, scaled\_n\_diagnoses, scaled\_n\_procedures, scaled\_n\_prescriptions, scaled\_comorbidity \\
\midrule
\multicolumn{2}{l}{\textbf{Structural Redundancy}} \\
& total\_icu\_los (redundant with los\_icu) \\
& comorbidity\_count (exact sum of icd\_* flags) \\
& received\_vasopressors (redundant with vasopressor\_duration\_hrs) \\
\midrule
\multicolumn{2}{l}{\textbf{EHR Frequency Artifacts}} \\
& All variables ending with \_n\_obs and \_n\_labs \\
& total\_output\_ml, n\_output\_events, n\_distinct\_cpt, n\_unique\_drugs, n\_unique\_drug\_classes \\
\bottomrule
\end{longtable}

\subsection{Character to Factor Conversion}

All remaining character variables are converted to factors:

\begin{verbatim}
mutate_if(is.character, as.factor)
\end{verbatim}

\subsection{Missing Data Handling}

Complete case analysis is applied:

\begin{verbatim}
dat <- na.omit(dat)
\end{verbatim}

\textbf{Note}: This may result in sample size reduction. The script reports row counts before and after NA omission.

\subsection{Final X/y and S/y\_S Extraction}
\label{sec:exclusions}

The target variables and predictor matrices are extracted in two steps:

\textbf{Step 1: Static Matrix (X, y)}
\begin{verbatim}
y <- dat$icu_mortality           # Binary outcome vector (N elements)
X <- as.data.frame(dat %>%
       select(-icu_mortality,
              -hadm_id))         # All remaining predictors (N x P)
rownames(X) <- valid_hadm_ids    # Row names = admission IDs
\end{verbatim}

\textbf{Step 2: Time-Varying Matrix (S\_mat, y\_S)}
\begin{verbatim}
S <- person_period_df %>%
  filter(hadm_id %in% valid_hadm_ids) %>%  # Match to cleaned cohort
  select(-any_of(exclude)) %>%             # Same exclusions as X
  # ... ethnicity recoding, na.omit() ...

y_S <- S$y_it                    # Discrete-time outcome (sum T_i elements)
S_mat <- as.data.frame(S %>%
           select(-y_it,
                  -icu_mortality,
                  -hadm_id))     # Time-varying predictors
\end{verbatim}

\textbf{Key difference}: \texttt{S\_mat} retains the \texttt{period} variable (integer 1--7) for baseline hazard approximation, while \texttt{X} has no period column (static analysis).

\subsection{Output Files}

Five output files are saved to \texttt{data/rds/}:

\begin{longtable}{L{4cm}L{10cm}}
\toprule
\textbf{File} & \textbf{Contents} \\
\midrule
\endhead
ready\_df.rds & Full static baseline matrix (before final exclusions) \\
longitudinal\_df.rds & Time-varying vital sign measurements \\
person\_period\_df.rds & Discrete-time survival format (before final exclusions) \\
\midrule
\textbf{X\_y.rds} & \textbf{List: \{X, y\}} --- Analysis-ready static predictors + outcome \\
\textbf{S\_y.rds} & \textbf{List: \{S, y\}} --- Analysis-ready time-varying predictors + discrete outcome \\
\bottomrule
\end{longtable}

\textbf{Recommended loading} (via \texttt{load.R}):

\begin{verbatim}
source("load.R")
# Automatically loads: X, y, S_mat, y_S
# Use X/y for BART, logistic regression
# Use S_mat/y_S for Polya-Gamma discrete-time models
\end{verbatim}

% ============================================================================
\section{Preprocessing Rationale}
% ============================================================================

\subsection{Age Capping}

MIMIC-III shifts dates of birth for patients $>$89 years (privacy protection), resulting in apparent ages $>$200. We cap at 90:

\begin{verbatim}
age = ifelse(age > 150, 90, age)
\end{verbatim}

\textbf{Justification}: Preserves information that patient is elderly without implausible values distorting regression models.

\subsection{Vital Sign Harmonization}

MIMIC-III has two charting systems (CareVue and MetaVision) with different item ID schemes. We map both to common variable names.

Temperature harmonization: CareVue recorded both Celsius and Fahrenheit. We convert Fahrenheit to Celsius: $(F - 32) \times 5/9$.

\subsection{Variable Scaling}

Continuous predictors are initially standardized (mean=0, SD=1) in \texttt{ready\_df}:

\begin{verbatim}
scaled_age = scale(age)[, 1]
\end{verbatim}

\textbf{However}: In Section 17, all \texttt{scaled\_*} variables are \textbf{removed} from the final \texttt{X} matrix to avoid multicollinearity with the original variables. If standardization is needed for modeling (e.g., Bayesian regression with weakly informative priors), it should be applied during model fitting.

\subsection{Missing Data Defaults}

\begin{itemize}[leftmargin=*, nosep]
    \item \textbf{Binary indicators}: Default to 0 (absence of code = condition not present)
    \item \textbf{Count variables}: Default to 0 (e.g., n\_diagnoses, n\_procedures)
    \item \textbf{DRG scores}: Default to 1 (lowest severity) --- conservative choice
    \item \textbf{Continuous vitals/labs}: Remain NA until factorization, then patients with missing values may be excluded via \texttt{na.omit()}
    \item \textbf{High-missingness variables}: Dropped entirely if $>25\%$ missing (see Section 5.2)
\end{itemize}

% ============================================================================
\section{Hierarchical Structure}
% ============================================================================

The data contains natural clustering that supports multi-level models:

\begin{enumerate}[leftmargin=*, nosep]
    \item \textbf{Patient level} (\texttt{subject\_id}): Some patients have multiple admissions
    \item \textbf{ICU level} (\texttt{first\_careunit}): 5 unit types (MICU, SICU, CCU, CSRU, TSICU)
    \item \textbf{Observation level}: For longitudinal\_df and person\_period\_df, repeated measurements/periods
\end{enumerate}

\textbf{Example hierarchical model}:
\[
h_{it} = h_0(t) \exp(X_i'\beta + Z_{i,t}'\gamma + u_i + v_{k[i]})
\]

where $u_i \sim N(0, \sigma^2_u)$ (patient frailty) and $v_k \sim N(0, \sigma^2_v)$ (ICU-level effect).

% ============================================================================
\section{Modeling Recommendations}
% ============================================================================

\subsection{Suggested Workflow}

\begin{enumerate}[leftmargin=*, nosep]
    \item \textbf{Exploratory}: Use \texttt{ready\_df} for descriptive statistics, univariate survival curves
    \item \textbf{BART modeling}: Use \texttt{X\_y.rds} directly --- contains pre-processed factors (e.g., *\_status, *\_group) for clinically interpretable tree-based analysis
    \item \textbf{Latent class discovery}: BART or Dirichlet process mixtures on \texttt{X} to identify heterogeneous patient strata
    \item \textbf{Discrete-time survival}: Bayesian logistic regression on \texttt{person\_period\_df} with Polya-Gamma augmentation
    \item \textbf{Joint modeling}: Link vital trajectories (\texttt{longitudinal\_df}) with survival via shared random effects
\end{enumerate}

\subsection{BART Implementation Notes}

When fitting BART with the \texttt{BART} package in R, use the pre-processed \texttt{X\_y.rds} file which has already:

\begin{enumerate}[leftmargin=*, nosep]
    \item \textbf{Excluded IDs and outcomes}: subject\_id, hadm\_id, icustay\_id, event\_status, hospital\_expire\_flag, time-to-event variables
    \item \textbf{Excluded timestamps}: All POSIXct columns
    \item \textbf{Excluded multicollinear variables}: Scaled versions, redundant aggregates
    \item \textbf{Converted to factors}: All categorical and discretized variables
    \item \textbf{Applied complete-case analysis}: NA rows removed
    \item \textbf{Recoded ethnicity}: Consolidated to 8 levels (0--7)
\end{enumerate}

\textbf{Key variables in final X matrix}:
\begin{itemize}[leftmargin=*, nosep]
    \item \textbf{Multi-level factors}: age\_group, los\_group, gcs\_status, hr\_status, rr\_status, spo2\_status, temp\_status, bun\_status, creat\_status, glucose\_status, hgb\_status, plt\_status, k\_status, na\_status, wbc\_status, comorbidity\_level, complexity\_status
    \item \textbf{Binary factors}: hr\_unstable, temp\_unstable, multiple\_icu\_stays, all icd\_* flags, intervention indicators
    \item \textbf{Categorical factors}: admission\_type, insurance, first\_careunit, first\_service, gender, ethnicity
    \item \textbf{Numeric}: age, los\_icu, n\_diagnoses, n\_procedures, n\_prescriptions, drg\_mortality\_score, drg\_severity\_score, n\_vasopressor\_doses, vasopressor\_duration\_hrs, n\_service\_changes, n\_transfers, n\_callouts, gcs\_mean\_val
\end{itemize}

Example R code for BART:
\begin{verbatim}
# Load pre-processed data
data <- readRDS("data/rds/X_y.rds")
X <- data$X
y <- data$y

# Fit logistic BART for binary outcome (icu_mortality)
library(BART)
bart_model <- lbart(x.train = X, y.train = y)
\end{verbatim}

\subsection{Prior Recommendations}

\textbf{Note}: The final \texttt{X} matrix primarily contains factor variables (both multi-level and binary). For any remaining numeric predictors (age, los\_icu, count variables), consider standardizing during model fitting.

For numeric predictors (when standardized):
\begin{align*}
\beta_j &\sim \text{Normal}(0, 2.5) \\
\sigma_u &\sim \text{Half-Cauchy}(0, 2.5)
\end{align*}

For period-specific intercepts in discrete-time models (baseline hazard):
\[
\alpha_t \sim \text{Normal}(0, 2) \quad \text{or random walk prior}
\]

For BART models, use default priors --- BART's regularization priors are well-calibrated for most applications.

% ============================================================================
\section{Data Quality Notes}
% ============================================================================

\begin{enumerate}[leftmargin=*, nosep]
    \item \textbf{NOTEEVENTS}: Empty in demo version (no clinical notes)
    \item \textbf{Age $>$89}: Capped at 90 due to MIMIC de-identification
    \item \textbf{ICD-9 codes}: Assigned retrospectively for billing; may not reflect status at ICU admission
    \item \textbf{Dual charting}: CareVue and MetaVision merged via item ID mapping
    \item \textbf{Lab missingness}: May be informative (stable patients get fewer labs)
    \item \textbf{LOCF limitations}: Artificially reduces variance; use missingness indicators as covariates
    \item \textbf{Demo vs. Full}: Demo has $\sim$100 patients; full has $\sim$46,000
\end{enumerate}

\end{document}
